\mode*
\section{Discussion}
\label{discussion}

\subsection{Implications}

If the hypotheses hold, then debugging skill is not a bag of tricks but a
particular case of a general learning ability: introducing the appropriate
patterns of variation for oneself.
Two consequences follow for teaching.

First, debugging can be taught \emph{as} the patterns of variation: instead of
teaching tools (debuggers, print statements) as such, we teach the students to
ask what to vary and what to hold invariant---contrast to open the dimension
where the bug lives, generalisation to isolate the critical aspect, fusion to
reconcile the new understanding with the old.
The worked example in \cref{debugging} is a template for such teaching.

Second, by the generative-learning effect
\autocites[pp.~227--228]{NCOL}{HolmqvistGustavssonWernberg2007}, teaching the
course content with
patterns of variation should over time improve the students' debugging even
without explicit debugging instruction---since debugging draws on the same
ability.
Since prgi26 is itself designed around the patterns of variation from the
companion paper \autocite{VTProgMisconceptions}, a first version of this
comparison is available already within the cohort: the students who study on
their own rather than from the course material (\cref{method}) do the same
labs without experiencing the material's patterns of variation.
A between-cohort comparison, against an instance taught without such
patterns, remains as a follow-up study.

\mode<presentation>{%
\begin{frame}
  \begin{block}{If the hypotheses hold}
    \begin{itemize}
      \item Debugging skill is a general learning ability, not a bag of
        tricks.
      \item Teach debugging as: what to vary, what to hold invariant.
      \item Teaching \emph{content} with patterns should improve debugging
        too (generative learning).
    \end{itemize}
  \end{block}
\end{frame}%
}

\subsection{Threats to validity}

The R-SPQ-2F is a self-report instrument; the deep-approach score measures
the student's reported approach to studying, not directly their behaviour in
programming.
This is partly by design---\cref{h:deep} is exactly the claim that the two
are connected---but a null result could mean either that the hypothesis is
wrong or that the instrument does not transfer to the programming context.
The original studies classified the approaches phenomenographically, from
task-specific accounts, and the questionnaire classification need not
coincide with that one; the agreement check against the interview-based
classification in the think-aloud subsample (\cref{think-aloud}) is
meant to catch this.

The learnlog data records the students' \emph{acts}---edits and runs---but
not the reasoning between them.
A student may generate a contrast mentally, without ever running it.
The think-aloud validation sessions (\cref{think-aloud}) are meant to bound this
gap, and the phenomenographic stance---acts express understanding---is
our theoretical justification for coding from acts, but the gap remains.

Confounds for \cref{h:success} include prior programming experience:
experienced students
will debug faster regardless of approach.
We measure prior programming knowledge with the start-of-course diagnostic
quiz and self-report (\cref{method}) and control for it.
The within-cohort comparison between material users and self-studiers
(\cref{h:generative}) is self-selected, not randomised: the students who skip the course material may
do so precisely because they are experienced, or because they already
generate their own variation---deep-approach students may be
over-represented among the self-studiers.
The R-SPQ-2F scores and the prior-experience measure let us control for this
statistically, but the comparison remains correlational.
Finally, the patterns of variation are our own analytic lens; the coding must
be validated with multiple independent coders and inter-rater reliability
reported.

\subsection{Conclusion}

We have argued, from variation theory, that debugging is self-directed
learning: a bug marks a critical aspect the programmer has not yet discerned,
and fixing it requires generating the patterns of variation---contrast,
generalisation, fusion---for oneself.
Since Marton connects the deep approach to learning to exactly this
self-generated variation, we hypothesise that debugging skill and the deep
approach are the same underlying ability.
A pilot with experienced programmers supports the first part: they could not
explain an unfamiliar phenomenon until they had generated all the patterns.
The planned study in prgi26---recording every edit--run cycle with learnlog,
measuring the approaches with R-SPQ-2F, and coding the episodes for patterns
of variation---will tell us whether the same holds for students, and whether
the deep approach predicts it.

\begin{frame}<presentation>
  \begin{idea}
    \begin{itemize}
      \item Debugging \(=\) learning without a teacher.
      \item Deep approach \(=\) generating one's own patterns of variation.
      \item Hypothesis: they are the same skill.
      \item prgi26 + learnlog + R-SPQ-2F will tell.
    \end{itemize}
  \end{idea}
\end{frame}
